\textbf{Пример 1}

Рассмотрим следующий вызов:

\begin{verbatim}
count_rectangles([[4,  8,  7,  5,  6],
                  [7,  4,  10, 3,  5],
                  [9,  7,  20, 14, 2],
                  [9,  14, 7,  3,  6],
                  [5,  7,  5,  2,  7],
                  [4,  5,  13, 5,  6]])
\end{verbatim}

\includegraphics{1.png}


Имеются $6$ допустимых областей, они перечислены ниже:
\begin{itemize}
\item $r_1 = r_2 = c_1 = c_2 = 1$
\item $r_1 = 1, r_2 = 2, c_1 = c_2 = 1$
\item $r_1 = r_2 = 1, c_1 = c_2 = 3$
\item $r_1 = r_2 = 4, c_1 = 2, c_2 = 3$
\item $r_1 = r_2 = 4, c_1 = c_2 = 3$
\item $r_1 = 3, r_2 = 4, c_1 = c_2 = 3$
\end{itemize}

Например, $r_1 = 1, r_2 = 2, c_1 = c_2 = 1$ является допустимой областью, так как
выполнены оба следующих условия:
\begin{itemize}
\item $a[1][1]=4$ строго меньше, чем $a[0][1]=8$, $a[3][1]=14$, $a[1][0]=7$, и $a[1][2]=10$.
\item $a[2][1]=7$ строго меньше, чем $a[0][1]=8$, $a[3][1]=14$, $a[2][0]=9$, и $a[2][2]=20$.
\end{itemize}
